\chapter{2001:William Knowles and Ryoji Noyori and K. Barry Sharpless}

\label{chap:chemistry2001}

The Nobel Prize in Chemistry for the year 2001 was awarded jointly to William Knowles and Ryoji Noyori and K. Barry Sharpless.

\textbf{Motivation:}

for their work on chirally catalysed hydrogenation reactions (one half, jointly, William Knowles and Ryoji Noyori) and for his work on chirally catalysed oxidation reactions (the other half, K. Barry Sharpless)

\section{William Knowles}

\subsection*{Early Life and Education}

William Knowles was born in 1917 in Taunton, Massachusetts, USA.

\subsection*{Career and Major Research}

Knowles studied at Harvard University (B.A., 1939) and earned his Ph.D. from Columbia University in 1942. He worked at the Monsanto Company, where he remained until his retirement in 1986.

\subsection*{Nobel Prize-winning Work}

The work for which William Knowles was awarded the Nobel Prize in Chemistry in 2001 was described by the Nobel Committee as: "for their work on chirally catalysed hydrogenation reactions".

He developed the first catalytic asymmetric hydrogenation, using a chiral rhodium catalyst, and applied it industrially to make L-DOPA, a drug used to treat Parkinson's disease.

\subsection*{Significance and Impact}

This was the first industrial example of asymmetric catalysis, producing a single enantiomer of a pharmaceutical.

\subsection*{Later Career and Honors}

Knowles retired from Monsanto in 1986. He shared the 2001 Nobel Prize in Chemistry, and he died in 2012.

\subsection*{Legacy}

His work showed that chiral catalysts could produce enantiopure compounds, and it opened the field of industrial asymmetric catalysis.

\subsection*{Selected Publications}

"Asymmetric Hydrogenation" (Accounts of Chemical Research, 1983)

\clearpage

\section{Ryoji Noyori}

\subsection*{Early Life and Education}

Ryoji Noyori was born in 1938 in Kobe, Japan.

\subsection*{Career and Major Research}

Noyori studied at Kyoto University (B.S., 1961; Ph.D., 1967). He was a professor at Nagoya University and later served as president of RIKEN from 2003 to 2015.

\subsection*{Nobel Prize-winning Work}

The work for which Ryoji Noyori was awarded the Nobel Prize in Chemistry in 2001 was described by the Nobel Committee as: "for their work on chirally catalysed hydrogenation reactions".

He developed the BINAP ligand and chiral ruthenium catalysts for highly enantioselective hydrogenation of a broad range of substrates.

\subsection*{Significance and Impact}

BINAP-based catalysis provided widely used methods for making chiral compounds, including pharmaceuticals and fragrances.

\subsection*{Later Career and Honors}

Noyori continued his research at Nagoya University and RIKEN. He shared the 2001 Nobel Prize in Chemistry.

\subsection*{Legacy}

His chiral molecular catalysts are among the most important tools in asymmetric synthesis.

\subsection*{Selected Publications}

"Asymmetric Catalysis by Architectural and Functional Molecular Engineering" (Angewandte Chemie International Edition, 2001)

\clearpage

\section{K. Barry Sharpless}

\subsection*{Early Life and Education}

K. Barry Sharpless was born in 1941 in Philadelphia, Pennsylvania, USA.

\subsection*{Career and Major Research}

Sharpless studied at Dartmouth College (B.A., 1963) and earned his Ph.D. from Stanford University in 1968. He held professorships at the Massachusetts Institute of Technology and Stanford University, and he joined the Scripps Research Institute in 1990.

\subsection*{Nobel Prize-winning Work}

The work for which K. Barry Sharpless was awarded the Nobel Prize in Chemistry in 2001 was recognized for his work on chirally catalysed oxidation reactions.

He developed catalytic asymmetric epoxidation and dihydroxylation, reliable methods for converting alkenes into chiral epoxides and diols with high enantioselectivity.

\subsection*{Significance and Impact}

The Sharpless oxidations became some of the most widely used reactions in asymmetric synthesis.

\subsection*{Later Career and Honors}

Sharpless continued his research at the Scripps Research Institute, and he received a second Nobel Prize in Chemistry in 2022 for the development of click chemistry.

\subsection*{Legacy}

His asymmetric oxidation reactions are standard tools for building chiral molecules, and his click chemistry concept transformed how chemists assemble molecules.

\subsection*{Selected Publications}

"The First Practical Method for Asymmetric Epoxidation" (Journal of the American Chemical Society, 1980)

\clearpage

\section*{Chapter Summary}

The 2001 Nobel Prize in Chemistry recognized catalytic methods for producing chiral molecules. Knowles and Noyori developed asymmetric hydrogenation, while Sharpless created asymmetric oxidation reactions.
